\part{Das levitische System}
Um ein gutes Verständnis von Jesus wirken als Opfer zu erlangen müssen wir zunächst einige Grundlagen im Alten Testament legen.
In diesem Teil werde ich für folgende Aussagen argumentieren:
\begin{enumerate}[(1)]
\item Kein Opfer im gesamten levitischen System kann als `substitution' des Opertieres an Stelle des Opfernden verstanden werden.
\item Bei keinem Opfer im levitischen System ist der `Tod' des Opfertieres zentral. Ganz im Gegenteil wird das Schlachten des Opfertieres rituell zu einem `nicht-töten' rekonzeptualisiert.
\item Es gibt kein Ritual oder Opfer im levitischen System, dass `Sünden vergeben' hat. Insbesondere sind die üblichen deutschen Übersetzungen `Sündenopfer', und `entsühnung' grundlegende Missverständnisse der Ziele und Limitierungen des levitischen Systems.
\end{enumerate}

\chapter{Grundlegende Konzepte des levitischen Systems}
Bevor wir die theologische Signifikanz des Opfersystems verstehen können, ist es nötig einige fundamentale Grundlagen des levitischen Systems zu verstehen.
Während wir dieses Grundverständnis erlangen ist es auch nötig, einige sehr übliche Wörter in den deutschen Übersetzungen des Alten Testaments aus unserem Sprachgebrauch zu entfernen.

\section{Heilig und Gemein}
Eine der wichtigsten Distinktionen, die das levitische System untermauert ist die Unterscheidung zwischen קֹדֶשׁ (qodeš, heilig) und חֹל (ḥol, gemein).
Heiliges ist (zu verschiednene Graden) abgesondert für einen speziellen Nutzen in Gottes Nähe.
Zunächst wertungsfrei sind gemeine Dinge oder Menschen solche, die nicht heilig sind.

Beispiele für heilige Dinge sind: das innere des Zeltes der Begegnung und des Tempels,
Gegenstände im Heiligtum,
die Kleider der Priester\bibfoot{Exodus}{40}{13},
die Priester selbst\bibfoot{Exodus}{40}{13},
das Fleisch verschiedener Opfer
\footnote{\bibverse{Leviticus}{10}{17}. Beachte, dass dieses Opferfleisch sogar hochheilig ist. Heiligkeit im levitischen System hat verschiedene Stufen.}
.

\section{Rein und Unrein}
Die zweite, und von Heiligkeit/Gemeinheit \textit{völlig} unabhängige Dimension ist die Unterscheidung zwischen טָחוֺר (ṭaḥor, rein) und טָמֵא (ṭame, unrein).
Im levitischen System gibt es verschiedene Arten von Reinheit, die es zu unterscheiden gilt.
Alle diese verschiednen Arten verwenden dieselben hebräischen Worte, sind aber in der Praxis leicht voneinander zu unterscheiden.
\begin{enumerate}[(a)]
	\item Jedes Ding und jeder Mensch ist zu einem bestimmten Zeitpunkt immer entweder \textit{rituell} rein oder \textit{rituell} unrein.
	Rituelle Unreinheit ist durch Berührung übertragbar\bibfoot{Numbers}{19}{22}.
	\item Manche Handlungen sind verboten, und dieses Verbot wird mit Reinheitssprache ausgedrückt.
	\footnote{In diese Kategorie fallen insbesondere die unreinen Tiere. Diese Tiere sind nicht grundsätzlich rituell unrein, denn Berührung eines unreinen Tieres überträgt keine rituelle Unreinheit. Ihr verzehr ist absolut verboten, und dieses Verbot drückt beispielsweise \bibverse{Leviticus}{11}{4} als unreinheit aus.}
	\item Moralisch verwerfliche Taten werden als moralische Unreinheit bezeichnet.
		\footnote{Wiederum ist offensichtlich, dass diese Art von Unreinheit unabhängig von ritueller Unreinheit ist. Das berühren eines rituell unreinen Menschen überträgt Unreinheit, das berühren eines moralisch unreinen Menschen (eines Sünders) überträgt keine rituelle (oder gar moralische) Unreinheit.}
\end{enumerate}
Es ist essentiell zu bemerken, dass heilig/unheilig und rein/unrein zwei verschiedene Dimensionen sind.
Ein Priester ist immer heilig (abgesondert für einen bestimmten Dienst für Gott), kann aber beispielsweise durch die Berührung einer Leiche unrein werden.
Das Beispiel aus \bibverse{Leviticus}{21}{1-3} ist instruktiv:
(a) Ein Priester kann sich Leichen-Unreinheit zuziehen wie jeder andere Mensch (b) er darf das im Normalfall nicht tun und (c) wenn es um Blutsverwandte geht ist dieses Verbot aufgehoben.
Wir sehen hieran die folgenden essentiellen Punkte, die wir über zwingend über rituelle Reinheit wissen müssen, bevor wir das Neue Testament lesen:
\begin{enumerate}[(1)]
	\item Sich rituelle Unreinheit zuzuziehen ist \textit{keine} Sünde.
	\item In manchen speziellen Situationen ist das zuziehen ritueller Unreinheit Sünde, und zwar \textit{aus göttlichem Befehl}. Diese Situation ist nicht der Normalzustand.
\end{enumerate}
Dass diese Punkte wahr sind sehen wir genauso auch an \bibverse{Leviticus}{15}{16-18}.
Es gibt keine Kinder ohne rituelle Unreinheit\footnote{Ganz zu schweigen von der Unreinheit der Frau nach der Geburt [\bibverse{Leviticus}{12}{}], oder der Unreinheit der Frau durch ihre Menstruation [\bibverse{Leviticus}{15}{19ff}.},
und es ist absolut normal, dass Menschen häufig rituell unrein werden.

Rituelle Unreinheit wird nur dann problematisch, wenn außerdem heilige Objekte im Spiel sind.
Eines der primären Probleme, dass das levitische System löst ist, dass heilige Objekte niemals mit Unreinheit in berührung kommen dürfen.
Wenn also eine Frau in den ersten sieben Tagen nach dem Beginn ihrer Menstruation (und dementsprechend unrein) ist, so darf sie nicht mit heiligen Gegenständen in Kontakt kommen, und nicht vom Fleisch von Dankopfern essen.
Es gibt keine weiteren Restriktionen oder Einschränkungen, die gewöhnliche rituelle Unreinheit nach sich zieht.
Eine Person die nicht vorhatte am Tempelkult in Jerusalem teilzunehmen hat also kaum einen Grund, viel über rituelle Unreinheit nachzudenken.
Ein rituelles Bad am Tag reicht aus, um den überwältigenden Großteil der rituellen Reinheitsgebote zu erfüllen.

Andere rituelle Unreinheiten können nicht durch ein einfaches Bad entfernt werden. Dazu zählt beispielsweise das berühren von Leichen\bibfoot{Numbers}{19}{}.
Die rituelle Unreinheit wird auch in solchen Fällen durch nicht-Opfer und Zeit entfernt, namentlich rituelle Waschungen, oder andere Riten (wie das besprengen der Person mit einer bestimmten Flüssigkeit).

Es gibt besonders große Fälle von ritueller Unreinheit, nämlich
Unreinheit durch (a) Aussatz\bibfoot{Leviticus}{14}{} (b) die Geburt eines Kindes\bibfoot{Leviticus}{12}{}, und (c) Blutfluss von Mann oder Frau\bibfoot{Leviticus}{15}{}.
Der entscheidende Unterschied zwischen großer und gewöhnlicher ritualler Unreinheit ist, dass große rituelle Unreinheit sich auf das Heiligtum überträgt, selbst ohne örtliche Proximität.
\Footcite[78f]{lamb-of-the-free}

Das bedeutet zum Beispiel, dass rituelle Unreinheit durch die Geburt eines Kindes auf das Heiligtum übertragen wird - das heilige kommt mit unreinem in Berührung, und für dieses Problem braucht es eine Lösung.
Wir kommen jetzt zu einem der wichtigsten Punkte, die wir über die Opfer im levitischen System verstehen müssen: Was ist `Entsühnung' und was bewirkt ein `Sündopfer'?

\section{\label{kipper}Dekontamination}
Die Lösung im levitischen System für Unreinheit, die auf das Heiligtum übertragen wurde ist die `Entsühnung', die ein `Sündopfer' bringen kann.

Das erste dieser Wörter übersetzt häufig das hebräische \kipper.
Die deutschen Glossen in einigen gängigen Übersetzungen sind:
`Sühnung erwirken'\footnote{\Cite{elb}, \Cite{meng}, \Cite{neü}, \Cite{ngü}, \Cite{slt}}, `Sühnung vollziehen'\Footcite{lut}, `Wiedergutmachung schaffen'\Footcite{nlb}, und `Versöhnung erwirken'\Footcite{eü}.

Alle diese Übersetzungen sind völlig falsch.\footnote{
Manchmal ist Klarheit nötig. Das Referenzwerk für die gesamte Diskussion um das levitische System ist \Cite{milgrom-leviticus-1}, siehe dort 255f für die Diskussion über \kipper.}

Das hebräische \kipper
\footnote{\lroot{כפר} im Piel}
hat als primäre Bedeutung `überdecken', wird aber insbesondere im levitischen System restriktiver verwendet: als `dekontaminieren' oder `reinigen' heiliger Gegenstände.\Footcite[4]{lamb-of-the-free}
Als anschauliche Begründung für diese Tatsache kann ein Vergleich des
`Sündopfers' für unwissentlich begangene Sünden des Volkes, \bibverse{Leviticus}{4}{20},
und des finale `Sündopfers' nach einer Geburt, \bibverse{Leviticus}{12}{7}, helfen.
Im ersten Fall lesen wir ``und der priester soll für sie kipper-wirken, und ihnen wird vergeben werden'', im zweiten ``und er soll kipper-wirken für sie, und sie wird rein sein''.
Wir sehen also, dass kipper sowohl zeitgleich mit `vergeben' als auch mit `rein werden' stehen kann.
Im ersten Fall gab es eine Sünde (und die damit verbundene verunreinigung des Heiligtums), im zweiten Fall gab es eine große rituelle Unreinheit (die sich auf das Heiligtum übertragen hat).
Im ersten Fall geschieht kipper und Vergebung, im zweiten Fall kipper und rein-Werden.
Das bedeutet, dass kipper eine Bedeutung haben muss, die völlig ohne Relation zu Sünde formulierbar ist (alternativ wäre das Gebähren eines Kindes eine Sünde).
Die Bedeutung die übrig bleibt\footnote{Also die maximale invariante (lexische) Semantik von \kipper, die diese beiden Stellen ohne einfügen eines homonym auflösen kann} ist ``Verunreinigung entfernen''.
Aus diesem Grund werde ich kipper durchgängig mit Dekontamination übersetzen und ab hier nur noch von `Entsühnung' sprechen, wenn das entfernen von Sünde tatsächlich ist was ich meine.

Eine weiteres fundamentales Missverständnis ist, dass das Blut von Opfern auf Menschen aufgetragen wird, um sie von Unreinheit und Schuld zu reinigen.
Im nächsten Kapitel werden wir alle Blutrituale ins Auge nehmen, bei denen Blut auf Menschen aufgetragen wird und werden sehen, dass kein Blut das als Teil eines Rituals Menschen berührt
\footnote{Menschen berühren soll. Es gibt Fälle, wo Opferblut von dekontaminierenden Opfern aus Versehen Dinge außer die heiligen Objekte berührt, die dekontaminiert werden sollen; das ist nicht beabsichtigt und problematisch.}
von einem dekontaminierenden Opfer stammt.
Aber das Verb kipper im gesamten Levitischen Text hat \textit{nie} eine Person als direktes Objekt.
Es wird `für eine Person dekontaminiert', aber das direkt Objekt dieser Dekontamination ist immer ein Objekt\footnote{oder elliptisch}, und niemals eine Person. Es wird `für eine Person ein Objekt dekontaminiert'.
Dieses Phänomen werden wir genauer untersuchen, wenn wir über das Reinigungsopfer (`Sündopfer') reden.


\chapter{Nicht dekontaminierende Opfer}
Was ist die primäre Aufgabe des levitischen Systems, und speziell des Opfersystems?
Wer in einer einfachen Version von PSA aufgewachsen ist könnte denken, dass die primäre Funktion der Opfer im Alten Testament das entfernen von Sünde und hindeuten auf den viel größeren stellvertretenden Opfertod des Messias ist.
Wir haben bereits im vergangenen Abschnitt festgehalten, dass kipper (das hinter den deutschen Übersetzungen `entsühnen' steht) nichts direkt mit Sünde oder Vergebung zu tun hat.
Was also ist dann die Funktion des Opfersystems?

Im großen ganzen sind das Opfersystem und das levitische System im ganzen eingesetzt, damit Israel mit ihrem Gott (der in ihrer Mitte wohnt) in Gemeinschaft stehen können.
So beginnt und endet jeder Tag mit einem regelmäßigen Brandopfer, das in \bibverse{Numbers}{28}{} eingesetzt wird, als "Gottes Speise", "ihm zum wohlgefälligen Geruch".
Dieses Brandopfer zieht die Aufmerksamkeit Gottes auf sich - ein Opfer, dass vollständig verbrannt wird und nur deshalb existiert, weil Gott es so wollte.
Neben den Brandopfern spielen die Heilsopfer eine große Rolle -
bei jedem der großen Jahresfeste werden Heilsopfer gebracht, um an die vergangenen Wohltaten Gottes zu erinnern, und auch um sein persönliches Dank an Gott auszudrücken hat ein Israelit ein Heilsopfer gebracht.
Die Heilsopfer sind speziell in der Hinsicht, dass ihr Fleisch unter anderem von der Person gegessen wird, die das Opfer bringt; es ist ein gemeinsames Mahl mit Gott:
ein drittel geht in Flammen vor Gott auf, ein drittel gehört den Priestern (die kein Land und deswegen keine eigenen Herden haben), ein drittel isst der Opfernde selbst.
Hier ist also in gröbsten Zügen die Funktion des Opfersystems:
(a) Brandopfer werden gebracht um Gott "anzuziehen", ihm als Wohlgeruch
(b) Heilsopfer werden dargebracht um Gott zu loben und ihm zu danken
(c) durch Unreinheit (moralisch und rituell) wird heiliger Ort verunreinigt; die dekontaminierenden Opfer sorgen dafür, dass diese Verunreinigung vom Heiligtum entfernt wird damit Gottes Heiligkeit weiterhin dort wohnen kann.


\section{Brandopfer und Heilsopfer}
- Beschreibung von Brandopfer und Heilsopfer
- Heilsopfer werden gegessen - nur nicht-dekontaminierende Opfer können gegessen werden
- Fußnote oder kurzer Satz: zum Teil hat olah wohl kipper-wirkung (cf. lev 1).
\section{Die Opfer zu Passah, Bundesschluss, und Bundeserneuerung}
- Brandopfer + Heilsopfer
- Erinnerung an Gottes Erlösung aus Ägypten
- Dankbarkeit für Gottes Bund
- Rededikation und neues Einladen Gottes nach Zeiten großer Sünde
\section{Blutrituale mit menschlichem Objekt}
- Metaphysischer Übergang
- Blut ist IMMER von nicht-dekontaminierenden Opfern

\chapter{Dekontaminierende Opfer}
- Funktion: entfernen von Unreinheit an heiligen Objekten
- Menschen werden IMMER UND AUSSCHLIEßLICH durch Waschung, Zeit oder andere Rituale gereinigt
\section{Das Reinigungsopfer}
- hat nichts mit Sünde zu tun
- Leben ist im Blut
- Tod des Tieres ist dafür da, um an Blut zu kommen
	- im vordergrund steht das Blutritual und was mit dem Fleisch des Tieres passiert
- Opfer im AT ist kein rituelles Schlachten; Schlachten ist eine notwendigkeit, um Blut und Fleisch zu haben
\section{Das Reparationsopfer}
- Unterart des Reinigungsopfers
- muss ich mir nochmal speziell anschauen
\section{Der Tag der Dekontamination}
- Ziel: Dekontamination des Heiligtums von innen nach außen
	- keine substitution; Opfer für das Volk, nicht Substitutionell an ihrer Stelle
- der Sündenbock kommt später
\section{Fallbeispiel: Reinigung nach der Geburt}
- Reinigung der Person durch Zeit und Waschung
- danach (Wenn die Person rein ist) hattat, um das Heiligtum zu kippern und "die Person wird rein"
- Mishna dazu zitieren; im prinzip aus lamb-of-the-free abschreiben

\chapter{Rituale, die keine Opfer sind}
% ich denke das kann man weglassen, weil die relevanten Punkte schon angesprochen sind
%\section{Reinigungsrituale für kleine rituelle Unreinheit}
%\section{Reinigungsrituale für große rituelle Unreinheit}
\section{\label{metaphysische-transformation}Der Bundesschluss und Blutreinigungsrituale}
- Verbindung zwischen Volk und Gott ausgedrückt als besprenklung
	- kein kipper, keine Vergebung
\section{Der Sündenbock}
- Sünde wird nicht kippert, sie wird übertragen und entfernt
- OPFER NEHMEN KEINE SÜNDE WEG
	- an dem einen Punkt, wo sünden aus dem Lager entfernt werden ist das KEIN OPFER
\section{Das erste Passah}
- kein opfer
- das muss ich mir nochmal genauer anschauen - was geht hier vor??
	- iwas mit "welchem Gott folgt ihr nach?", weil das Lamm ein Idol von Atum ist?? steckt da was dahinter?

%% ggf. auch direkt oben eingebaut
\chapter{Was Opfer nicht sind}
\section{Substitution}
das Handauflegungsritual
\section{\label{opfer-als-nicht-tod}Tod, Strafe, und Leiden}
Rituelle Rekonzeption des Todes
Der Tod des Tieres ist ein Mittel Zum zweck: um an Blut und Fleisch zu kommen

